\begin{name}
	{\tenchude}
	{TOÁN 10}
	{LỚP TOÁN THẦY PHÁT}
	{{Ai cũng sẽ phải trưởng thành dù muốn hay không.\\ Vậy hãy chọn cách mạnh mẽ mà trưởng thành.}}
\end{name}
\TN
\begin{ex}%[Dự án Form mới 12/4/6-HK1 Lớp 10-11]%[VN-MT-6, Võ NGuyên Thạch]%[THPT TÂN TÚC-TP HCM]%[0D6N3-3]
Trung vị của mẫu số liệu $4$; $5$; $8$; $13$; $5$; $17$; $9$ là
\choice
{$9$}
{$5$}
{$13$}
{\True $8$}
\loigiai{
Sắp xếp mẫu số liệu thành dãy không giảm $4$; $5$; $5$; $8$; $9$; $13$; $17$.\\
Trung vị của mẫu số liệu là $M_e=8$.
}
\end{ex}

\begin{ex}%[0-HK1-CT-9-TranPhu-HoChiMinh-2324]%[VN-MT-6,Nguyễn Văn Minh Hiếu]%[0D6N4-2]
Điểm thi học kì I của một học sinh lớp 10 như sau: $9$; $9$; $7$; $8$; $9$; $7$; $10$; $8$; $8$. Khoảng biến thiên của mẫu số liệu là
\choice
{$1$}
{$2$}
{\True$3$}
{$4$}
\loigiai{
Xếp lại mẫu số liệu $7; 7; 8; 8; 8; 9; 9; 9; 9; 10$.\\
Khoảng biến thiên của mẫu số liệu là $10-7=3$.
}
\end{ex}

\begin{ex}%[0D8H1-3]
Từ tập $\{1;2;3;4;5;6\}$ lập được bao nhiêu số tự nhiên có nhiều nhất hai chữ số?
\choice
{$30$}
{\True $42$}
{$36$}
{$6$}
\loigiai
{
\begin{itemize}
\item Số một chữ số có $6$ số.
\item Số có hai chữ số có $6^2=36$ số.
\end{itemize}
Vậy có tất cả $36+6=42$ số thỏa mãn yêu cầu bài toán.
}
\end{ex}

\begin{ex}%[0D8H2-3]
Có bao nhiêu số có ba chữ số dạng $\overline{abc}$ với $a, b, c\in\left\{{0;1;2;3;4;5;6}\right\}$ sao cho $ a < b < c$.
\choice
{$ 30$}
{\True $ 20$}
{$ 120$}
{$ 40$}
\loigiai{
Số các số tự nhiên thỏa mãn bài ra bằng số các tổ hợp chập $ 3$ của $ 6$ phần tử thuộc tập hợp $\left\{{1,2,3,4,5,6}\right\}$.\\
Vậy có $ \mathrm{C}_6^3=20$ số.}
\end{ex}

\begin{ex}%[BG - 10 New - 4in1,Phan Chiến Thắng]%[0D8H2-4]
Số cách sắp xếp $6$ nam sinh và $4$ nữ sinh vào một dãy ghế hàng ngang có $10$ chỗ ngồi là
\choice
{$6!-4!$}
{$6!4!$}
{\True $10!$}
{$6!+4!$}
\loigiai{
Số cách sắp xếp $6$ nam sinh và $4$ nữ sinh vào một dãy ghế hàng ngang có $10$ chỗ là số hoán vị của $10$ phần tử nên có $10!$ cách.
}
\end{ex}

\begin{ex}%[0D8H2-6]
Cho $10$ điểm phân biệt và không có ba điểm nào thẳng hàng. Hỏi có bao nhiêu tam giác có $3$ đỉnh là $3$ trong $10$ điểm đã cho?
\choice
{$720$}
{$27$}
{\True $120$}
{$1\,000$}
\loigiai{
Mỗi cách chọn $3$ trong $10$ điểm đã cho để được $3$ đỉnh của một tam giác là một tổ hợp chập $3$ của $10$ phần tử.\\
Vậy có $\mathrm{C}_{10}^3 = 120$ tam giác.
}
\end{ex}

\begin{ex}%[0D8H3-2]
Khai triển nhị thức $(a-2b)^5$ thành tổng các đơn thức, ta được kết quả là
\choice
{$a^5+ 10a^4b- 40a^3b^2+ 80a^2b^3- 80ab^4+ 32b^5$}
{$a^5- 10a^4b+ 40a^3b^2- 80a^2b^3+ 40ab^4- b^5$}
{$a^5- 5a^4b+ 10a^3b^2- 10a^2b^3+ 5ab^4- b^5$}
{\True $a^5- 10a^4b+ 40a^3b^2- 80a^2b^3+ 80ab^4- 32b^5$}
\loigiai{
Ta có \begin{align*}
(a-2b)^5 = &\,a^5+ 5a^4(-2b)+ 10a^3(-2b)^2+ 10a^2(-2b)^3+ 5a(-2b)^4+ (-2b)^5\\
= &\, a^5- 10a^4b+ 40a^3b^2- 80a^2b^3+ 80ab^4- 32b^5.
\end{align*}  }
\end{ex}

\begin{ex}%[0H9H1-2]
Trong mặt phẳng $Oxy$, cho $4$ điểm $M(1;-1)$, $N(5;-3)$, $P(0;4)$ và $G(2;0)$. Đẳng thức nào dưới đây là \textbf{sai}?
\choice
{$\overrightarrow{ON}+\overrightarrow{OM}+\overrightarrow{OP}=3\overrightarrow{OG}$}
{$\overrightarrow{GN}+\overrightarrow{GM}+\overrightarrow{GP}=\overrightarrow{0}$}
{\True $\overrightarrow{GN}+\overrightarrow{GM}=\overrightarrow{GP}$}
{$\overrightarrow{NG}+\overrightarrow{MG}+\overrightarrow{PG}=\overrightarrow{0}$}
\loigiai{
Ta có $\overrightarrow{GN}=(3;-3)$, $\overrightarrow{GM}=(-1;-1)$, $\overrightarrow{GP}=(-2;4)$. \\
Do đó $\overrightarrow{GN}+\overrightarrow{GM}=(3+(-1);-3+(-1))=(2;-4) $ và $\overrightarrow{GP}=(-2;4)$.\\
Vậy $\overrightarrow{GN}+\overrightarrow{GM} \neq \overrightarrow{GP}$ nên $\overrightarrow{GN}+\overrightarrow{GM}=\overrightarrow{GP}$ là đẳng thức sai.}
\end{ex}

\begin{ex}%[0H9H3-2]
Trong mặt phẳng tọa độ $O x y$, cho hai điểm $A(1 ;-2)$ và $B(3 ; 2)$. Phương trình tổng quát của đường thẳng $A B$ là
\choice
{$2 x+4 y+6=0$}
{$2 x-y+4=0$}
{$x+2 y-10=0$}
{\True $2 x-y-4=0$}
\loigiai{
Ta có $\overrightarrow{AB}=(2;4)$ là véc-tơ chỉ phương của $AB$ suy ra véc-tơ pháp tuyến của $AB$ là $\overrightarrow{n}=(4;-2)$.
Phương trình tổng quát của đường thẳng $A B$ đi qua $A(1;-2)$ và nhận $\overrightarrow{n}=(4;-2)$ làm véc-tơ pháp tuyến là
\[4(x-1)-2(y+2)=0\Leftrightarrow 4x-2y-8=0 \Leftrightarrow 2x-y-4=0.\]
}
\end{ex}

\begin{ex}%[Dự án Đề GK2 PNL Cấu Trúc 2025]%[Đào-V- Thuỷ]%[0H9H4-2]
Trong mặt phẳng $Oxy$, cho hai điểm $A(-1;-3)$, $B(-3;5)$. Phương trình đường tròn có đường kính $AB$ là
\choice
{\True $(x+2)^2+(y-1)^2=17$}
{$(x+2)^2+(y-1)^2=\sqrt{17}$}
{$(x+1)^2+(y-4)^2=\sqrt{68}$}
{$(x+1)^2+(y+3)^2=68$}
\loigiai
{
Đường tròn đường kính $AB$ có tâm $I$ là trung điểm của $AB$. Suy ra $I(-2; 1)$.\\
Khi đó bán kính $R= IA= \sqrt{1^2+(-4)^2}= \sqrt{17}$.\\
Vậy phương trình đường tròn cần tìm là $(x+2)^2+(y-1)^2=17$.
}
\end{ex}

\begin{ex}%[0H9H5-2]%[Dự án đề kiểm tra Toán khối 11 HKII NH23-24-Dot15-PhucChuong]%[THPT Nhu Van Lang - Hai Phong]
Dạng chính tắc của Elip là
\choice
{\True$\dfrac{x^2}{a^2}+\dfrac{y^2}{b^2}=1$}
{$\dfrac{x^2}{a^2}-\dfrac{y^2}{b^2}=-1$}
{$\dfrac{x^2}{a^2}+\dfrac{y^2}{b}=1$}
{$\dfrac{x^2}{a^2}-\dfrac{y^2}{b^2}=1$}
\loigiai{
Dạng chính tắc của Elip là $\dfrac{x^2}{a^2}+\dfrac{y^2}{b^2}=1$.
}
\end{ex}

\begin{ex}%[0H9H5-4]%[Dự án đề kiểm tra Toán 10 CHKI NH23-24- lamnguyen]%[THPT - Đồng nai]
Cho hypebol có phương trình: $\dfrac{x^2}{16}-\dfrac{y^2}{9}=1$. Tiêu cự của hypebol là
\choice
{$\sqrt{7}$}
{$2\sqrt{7}$}
{$5$}
{\True $10$}
\loigiai{
Ta có $a^2=16$, $b^2=9$\\
$c^2=a^2+b^2=16+9=25\Rightarrow c=5\Rightarrow 2c=10$.}
\end{ex}
\TNTF
\begin{ex}%[BG - 10 New - 4in1,Phan Chiến Thắng]%[0D8H2-4]
Sắp xếp năm bạn học sinh An, Bình, Chi, Dũng, Lệ vào một chiếc ghế dài có $5$ chỗ ngồi.
\choiceTF[t]
{\True Số cách sắp xếp $n$ người ngồi vào dãy ghế dài $n$ chỗ ngồi là $n!$}
{\True Số cách xếp $5$ bạn vào $5$ chỗ trên ghế dài là $120$}
{Số cách xếp sao cho bạn An và bạn Dũng luôn ngồi cạnh nhau là $24$}
{\True Số cách sắp xếp sao cho bạn An và Dũng không ngồi cạnh nhau là $72$}
\loigiai{
\begin{itemchoice}
\itemch Đúng. Số cách sắp xếp $n$ người ngồi vào dãy ghế dài $n$ chỗ ngồi là $n!$
\itemch Đúng. Số cách xếp $5$ bạn vào $5$ chỗ trên ghế dài là số hoán vị của $5$ phần tử nên có $5!=120$ cách.
\itemch Sai. Số cách xếp sao cho bạn An và bạn Dũng luôn ngồi cạnh nhau là $2\cdot 4!=48$ cách (An và Dũng ngồi cạnh nhau xem như $1$ bạn; xếp $4$ bạn vào $4$ chỗ có $4!$ cách; cách xếp An và Dũng ngồi cạnh nhau là $2!=2$).
\itemch Đúng. Số cách sắp xếp sao cho bạn An và bạn Dũng không ngồi cạnh nhau là $120-48=72$ cách.
\end{itemchoice}
}
\end{ex}

\begin{ex}%[Dự án 18 - TeamTeXHoa - Phạm Quốc Toàn]%[0H9H4-5]
Trong hệ trục tọa độ $Oxy$, cho đường tròn $\left(C \right)$ tâm $I \left(1;2 \right)$ và cắt đường thẳng \\
$\Delta \colon 3x+4y-6=0$ tại hai điểm $A$, $B$ sao cho $S_{IAB}=4$. Trong mỗi ý a), b), c), d) sau, hãy chọn đúng hoặc sai.
\choiceTF
{\True Khoảng cách từ tâm $I$ đến đường thẳng $\Delta$ bằng $1$}
{Bán kính đường tròn $\left(C \right)$ nhỏ hơn $4$}
{Phương trình đường tròn $\left(C \right) \colon x^2+y^2-2x-4y+12=0$}
{Điểm $O$ nằm trên đường tròn $\left(C \right)$}
\loigiai{
\begin{center}
\begin{tikzpicture}[>=stealth, x=1.0 cm, y=1.0 cm]
% \draw[->] (-4,0) --(0,0) node[below left]{$O$} -- (4,0) node[below]{$x$};
% \draw[->](0,-4)--(0, 4) node[right]{$y$};
\draw(0,0) circle (3); %(0,0) là tọa độ tâm, 1 là bán kính
\draw(-4,2) -- (4,2) (0, 0) -- (0, 2) (0, 0) -- (-2.24, 2) (0, 0) -- (2.24, 2);
\draw (0,0) circle (1pt) node[below]{$I$};
\draw (0,2) circle (1pt) node[above]{$H$};
\draw (3.5,2) circle node[below]{$\Delta$};
\draw (-2.24,2) circle (1pt) node[above left]{$A$};
\draw (2.24,2) circle (1pt) node[above right]{$B$};
\end{tikzpicture}
\end{center}
\begin{itemchoice}
\itemch Gọi $H$ là trung điểm của đoạn thẳng $AB \Rightarrow IH \perp AB$. Ta có $IH=d_{\left(I, \Delta\right)}=\dfrac{\left| 3+8-6\right|}{\sqrt{3^2+4^2}}=1$.
\itemch $IH \perp AB \Rightarrow S_{IAB}=\dfrac{1}{2} AB\cdot IH \Leftrightarrow 4=\dfrac{1}{2}AB\cdot 1\Leftrightarrow AB=8$ \\
$\Rightarrow AH=4$ và $R=\sqrt{AH^2+IH^2}=\sqrt{17}$ $\Rightarrow R>4$.
\itemch Phương trình đường tròn $\left(C \right)$ tâm $I \left(1;2 \right)$, bán kính $R=\sqrt{17}$ là $\left(C \right) \colon \left(x-1 \right)^2+ \left(y-2 \right)^2=17$ hay $\left(C \right) \colon x^2+y^2-2x-4y-12=0$.
\itemch Ta có $\overrightarrow{IO}= \left(-1;-2 \right) \Rightarrow IO=\sqrt{\left(-1 \right)^2+ \left(-2 \right)^2}=\sqrt{5}<R=\sqrt{17}$. Suy ra $O$ nằm trong $\left(C \right)$.
\end{itemchoice}
}
\end{ex}
\TNSA
\begin{ex}%[Dự án đề theo cấu trúc mới 2025-NCT]%[0D0V2-5]
Cho đa giác đều gồm $2n$ đỉnh. Chọn ngẫu nhiên ba đỉnh trong số $2n$ đỉnh của đa giác, xác suất ba đỉnh được chọn tạo thành một tam giác vuông là $0{,}2$. Tìm $ n$, biết $ n$ là số nguyên dương và $ n\ge 2$.
\shortans{$8$}
\loigiai{
Do chọn ngẫu nhiên $3$ đỉnh từ $2n$ đỉnh, nên số phần tử của không gian mẫu là $ n(\Omega)=\mathrm{C}_{2n}^3$.\\
Ba đỉnh được chọn tạo thành tam giác vuông khi và chỉ khi có hai đỉnh trong ba đỉnh là hai đầu mút của một đường kính của đường tròn ngoại tiếp đa giác, và đỉnh còn lại là một trong số $( 2n-2 )$ đỉnh còn lại của đa giác.\\
Số cách chọn một đường kính (mà hai đâu mút là hai đỉnh trong số $2n$ đỉnh của đa giác) là $ n$.\\ Suy ra số cách chọn $3$ đỉnh trong số $2n$ đỉnh của đa giác để chúng tạo thành một tam giác vuông là $ n( 2n-2 )$ .\\
Vì xác suất ba đỉnh được chọn tạo thành một tam giác vuông là $0{,}2$, nên \[\dfrac{n\left(2n-2\right)}{\mathrm{C} _{2n}^3}=0{,}2\Leftrightarrow\dfrac{6n\left(2n-2\right)}{2n(2n-1)(2n-2)}=\dfrac 15\Leftrightarrow n=8.\]
}
\end{ex}

\begin{ex}%[0D8H2-5]
Một hộp đựng $11$ tấm thẻ được đánh từ $1$ đến $11$. Hỏi có bao nhiêu cách lấy ra $6$ thẻ để tổng số ghi trên $6$ thẻ đó là một số lẻ.
\shortans{$236$}
\loigiai{
Ta ký hiệu: thẻ ghi số lẻ là thẻ lẻ; thẻ ghi số chẵn là thẻ chẵn.\\
Tổng số ghi trên 6 thẻ lấy ra là một số lẻ thì số thẻ lẻ lấy ra phải là một số lẻ.\\ Ta có các trường hợp sau:
\begin{itemize}
\item TH1: $1$ thẻ lẻ, $5$ thẻ chẵn, suy ra $\mathrm{C}_6^1 \cdot \mathrm{C}_5^5=6$ cách.
\item TH2: $3$ thẻ lẻ, $3$ thẻ chẵn, suy ra $\mathrm{C}_6^3 \cdot \mathrm{C}_5^3=200$ cách.
\item TH3: $5$ thẻ lẻ, $1$ thẻ chẵn, suy ra $\mathrm{C}_6^5 \cdot \mathrm{C}_5^1=30$ cách.
\end{itemize}
Vậy có: $6+200+30=236$ cách.
}
\end{ex}

\begin{ex}%[0D8H3-4]
Tính tổng các hệ số của số hạng chứa $x^n$ với $n\in\mathbb{N}$ trong khai triển nhị thức Newton $\left(x^2+\dfrac{2}{x}\right)^5$, với $x\ne0$.
\shortans{$131$}
\loigiai{Ta có
\begin{eqnarray*}
\left(x^2+\dfrac{2}{x}\right)^5&=&\mathrm{C}_5^0\left(x^2\right)^5+\mathrm{C}_5^1\left(x^2\right)^4\left(\dfrac{2}{x}\right)+\mathrm{C}_5^2\left(x^2\right)^3\left(\dfrac{2}{x}\right)^2+\mathrm{C}_5^3\left(x^2\right)^2\left(\dfrac{2}{x}\right)^3+\mathrm{C}_5^4\left(x^2\right)\left(\dfrac{2}{x}\right)^4+\mathrm{C}_5^5\left(\dfrac{2}{x}\right)^5\\
&=&x^{10}+10x^7+40x^4+80x+\dfrac{80}{x^2}+\dfrac{32}{x^5}.
\end{eqnarray*}
Vậy tổng các hệ số cần tìm là $T=1+10+40+80=131$.}
\end{ex}

\begin{ex}%[0H9V4-7]%[Dự án  Toán thực tế]%[Huỳnh Xuân Tín]
\immini[thm]{Ném đĩa là một môn thể thao thi đấu trong Thế vận hội Olympic mùa hè. Khi thực hiện cú ném, vận động viên thường quay lưng lại với hướng ném, sau đó xoay ngược chiều kim đồng hồ một vòng rưỡi của đường tròn để lấy đà rồi thả tay ra khỏi đĩa. Giả sử đĩa chuyển động trên một đường tròn tâm $I\left(0;\dfrac{3}{2}\right)$ bán kính $0{,}8$ trong mặt phẳng tọa độ $Oxy$ (đơn vị trên hai trục là mét). Đến điểm $M\left(\dfrac{\sqrt{39}}{10};2\right)$, đĩa được ném đi (hình bên). Trong những giây đầu tiên ngay sau khi được ném đi, quỹ đạo chuyển động của chiếc đĩa là phương trình có dạng $ax+by+\dfrac{261}{100}=0$. Tính $10a+2b$ (làm tròn đến hàng phần trăm)
}
{
\begin{tikzpicture}[>=stealth,line join=round,line cap=round,font=\footnotesize,scale=.7,xscale=0.8, yscale=0.8]
\draw[->] (-3,0)--(3,0) node[below]{$x$};
\draw[->] (0,-1)--(0,7) node[left]{$y$};
\coordinate (O) at (0,0);
\coordinate (I) at (0,1.5);
\coordinate (M) at (0.624,2);
\coordinate (a) at ($(M)!4!-90:(I)$);
\draw (I) circle(0.8 cm) (M)--(a);
\draw[fill=black]
(I) circle(1pt) node[left,scale=0.8]{$I$}
(0,2) circle(1pt) node[left,scale=0.8]{$2$}
(M) circle(1pt) node[right,scale=0.8]{$M$}
(a) circle(1pt)
($(a)+(-0.3,0.2)$) circle(1pt)
($(a)+(-0.6,0.4)$) circle(1pt)
($(a)+(-0.8,0.5)$) circle(1pt)
;
\end{tikzpicture}
}
\shortans{$7{,}24$}
\loigiai{
Xét đường tròn $(C)$ tâm $I\left(0;\dfrac{3}{2}\right)$, bán kính $R=0{,}8$.\\
Ta có $IM=\sqrt{\left(\dfrac{\sqrt{39}}{10}-0\right)^2+\left(2-\dfrac{3}{2}\right)^2}=0{,}8$.\\
Suy ra điểm $M$ thuộc đường tròn tâm $I$, bán kính $R=0{,}8$.\\
Trong những giây đầu tiên sau khi ném đi, quỹ đạo chuyển động của chiếc đĩa là một đường thẳng có phương trình là phương trình tiếp tuyến của đường tròn tại điểm $M\left(\dfrac{\sqrt{39}}{10};2\right)$
\[\left(\dfrac{\sqrt{39}}{10}-0\right)\left(x-\dfrac{\sqrt{39}}{10}\right)+\left(2-\dfrac{3}{2}\right)(y-2)=0\Leftrightarrow\dfrac{\sqrt{39}}{10}x+\dfrac{1}{2}y+\dfrac{261}{100}=0.\]
Ta có $a=\dfrac{\sqrt{39}}{10}$ và $b=\dfrac{1}{2}$ nên $10a+2b=10\cdot \dfrac{\sqrt{39}}{10}+2\cdot \dfrac{1}{2}\approx 7{,}24$.
}
\end{ex}
\TL
\begin{ex}%[0D8H3-2]%[Dự án đề kiểm tra Toán 10 HK2 NH23-24- Don Lee]%[THPT Gia Định - Tp Hồ Chí Minh]
Tìm hệ số của số hạng chứa $x^3$ trong khai triển $(2x-3)(3x+1)^4$.
\loigiai{
Ta có
\allowdisplaybreaks
$\begin{aligned}[t]
(2x-3)(3x+1)^4&=(2x-3)\left(81x^4+108x^3+54x^2+12x+1\right)\\
&=162x^5-27x^4-216x^3-138x^2-34x-3.
\end{aligned}$\\
Vậy hệ số của số hạng chứa $x^3$ trong khai triển là $-216$.
}
\end{ex}

\begin{ex}%[24-25 giảng K10 - K11, Phan Quốc Trí]%[0D8V2-4]
Đội văn nghệ của một nhà trường gồm $4$ học sinh lớp $12$A, $3$ học sinh lớp $12$B và $2$ học sinh lớp $12$C.Tìm số cách chọn $5$ học sinh từ đội văn nghệ để biểu diễn trong lễ bế giảng sao cho lớp nào cũng có học sinh và có ít nhất $2$ học sinh lớp $12$A.

\loigiai{
Từ giả thiết ta thấy có các trường hợp sau
\begin{center}
\begin{tabular}{|c|c|c|c|}
\hline
Số học sinh $12$A & Số học sinh $12$B & Số học sinh $12$C & \text{Số cách chọn}\\
\hline
$2$ & $2$ & $1$ & $\mathrm{C}^2_4 \times \mathrm{C}^2_3 \times \mathrm{C}^1_2 = 36$\\
\hline
$2$ & $1$ & $2$ & $\mathrm{C}^2_4 \times \mathrm{C}^1_3 \times \mathrm{C}^2_2 = 18$\\
\hline
$3$ & $1$ & $1$ & $\mathrm{C}^3_4 \times \mathrm{C}^1_3 \times \mathrm{C}^1_2 = 24$\\
\hline
\end{tabular}
\end{center}
Vậy có tất cả $36+18+24=78$ cách chọn thoả mãn yêu cầu bài toán.
}
\end{ex}

\begin{ex}%[0H9V4-5]%[Dự án đề kiểm tra Toán 10 HK2 NH23-24-BCTuan]%[THPT Nguyễn Thị Minh Khai-TPHCM]
Viết phương trình đường tròn đi qua $A(-8;-1)$ và tiếp xúc với đường thẳng $d\colon 3x+y+9=0$ tại $B(-4;3)$.
\loigiai{
Gọi $I$ là tâm của đường tròn, khi đó $IB\perp d$.\\
Phương trình của đường thẳng $IB$ đi qua $B(-4;3)$ và vuông góc với $d$ là $x-3y+13=0$.\\
Vì $IA=IB$ nên $I$ thuộc đường trung trực $\Delta$ của $AB$.\\
Phương trình đường thẳng $\Delta$ qua $M(-6;1)$ là trung điểm của $AB$ và vuông góc với $AB$ là $x+y+5=0$.\\
Tọa độ $I$ là nghiệm của hệ $\heva{& x-3y+13=0 \\ & x+y+5=0}\Leftrightarrow \heva{& x=-7 \\ & y=2.}\Rightarrow I(-7;2)$.\\
Khi đó $R=IA=\sqrt{10}$.\\
Vậy phương trình của đường tròn là $(x+7)^2+(y-2)^2=10$.
}
\end{ex}

